\section{Introduction}
\label{sec:intro}
\IEEEPARstart{S}{patio-temporal} forecasting predicts future states in graph-structured systems (e.g., transportation networks, power grids) where temporal dynamics and spatial interactions evolve jointly. Node-level trajectories are rarely self-driven; rather, they are shaped by historical conditions, future external drivers, and domain-specific feasibility constraints. As illustrated in Fig.~\ref{fig:multi_scale_processes}, complex network signals comprise shared macro-trends driven by exogenous forcing and high-frequency local variations. Consequently, a practical forecasting model must learn continuous temporal representations, infer dynamic graph relations, and integrate future exogenous variables (ExG) without data leakage, all while maintaining operational compatibility.

\begin{figure}[htbp]
\centering
\begin{tikzpicture}[font=\sffamily\tiny, >=Stealth]
    % Coordinates and grid
    \draw[help lines, gray!10, step=0.5] (0,-1.2) grid (6.2,1.5);
    \draw[->, thick, gray!60] (0,0) -- (6.5,0) node[right, black] {Time};
    \draw[->, thick, gray!60] (0,-1.2) -- (0,1.8) node[above, black] {System State};
    
    % Macro trend (Blue)
    \draw[very thick, blue!70!black, smooth, opacity=0.8] 
        plot coordinates {(0.1,-0.55) (0.9,-0.25) (1.8,0.15) (2.8,0.48) (3.8,0.32) (4.8,0.75) (6.0,0.98)}
        node[right=2pt, yshift=4pt, font=\sffamily\tiny\bfseries, blue!70!black] {Macro Trend};
    
    % Local variation (Orange)
    \draw[thick, orange!90!black, smooth] 
        plot coordinates {(0.1,-0.45) (0.35,-0.2) (0.62,-0.65) (0.9,-0.18) (1.15,-0.42) (1.45,0.02) (1.8,-0.08) (2.1,0.38) (2.45,0.02) (2.8,0.55) (3.1,0.2) (3.45,0.72) (3.8,0.25) (4.1,0.64) (4.5,0.38) (4.8,0.88) (5.2,0.58) (5.6,1.12) (6.0,0.9)}
        node[right=2pt, yshift=-6pt, font=\sffamily\tiny\bfseries, orange!90!black] {Node-specific Noise};

    % Callouts
    \node[draw=blue!60, fill=blue!5, rounded corners=2pt, align=center, inner sep=2pt] at (1.4,1.3) 
        {\textbf{Global Forcing}\\Shared low-frequency ExG};
    \draw[->, blue!60, shorten >=2pt] (1.4,1.0) -- (2.5,0.4);

    \node[draw=orange!60, fill=orange!5, rounded corners=2pt, align=center, inner sep=2pt] at (4.5,-0.8) 
        {\textbf{Local Perturbations}\\High-frequency spatial jitter};
    \draw[->, orange!60, shorten >=2pt] (4.5,-0.5) -- (5.0,0.8);

    % Legend-like markers (Moved higher to avoid callout)
    \fill[blue!70!black] (0.2,1.85) circle (1.5pt) node[right, black] {Shared Pattern};
    \fill[orange!90!black] (2.2,1.85) circle (1.5pt) node[right, black] {Local Variation};

\end{tikzpicture}
\caption{Multi-scale spatio-temporal processes. General network signals are composed of shared macro-trends driven by exogenous forcing and high-frequency variations unique to individual nodes. PhyST-Net separates these components through dual-branch graph reasoning.}
\label{fig:multi_scale_processes}
\end{figure}

\subsection{The Forecasting Trilemma in Spatio-Temporal Networks}
\label{sec:trilemma}
Despite recent architectural advancements, modern spatio-temporal forecasting frameworks remain hampered by a ``Forecasting Trilemma''---an inherent trade-off between Model Capacity, Domain Consistency, and Spatio-Temporal Complexity. 

\subsubsection{Model Capacity vs. High-Frequency Noise}
\label{sec:capacity}
Capacity denotes architectural expressive power. While modern Transformers excel at sequence modeling, they are highly sensitive to high-frequency stochastic noise in raw spatio-temporal data. Furthermore, rigid temporal patching exacerbates this by artificially splitting continuous physical regimes. As highlighted by recent trends in complexity-balanced sequence modeling, uniform representational capacity across an entire time horizon is computationally inefficient. The challenge is designing an architecture capable of capturing macro-trends while employing adaptive temporal allocation to filter node-specific noise.

\subsubsection{Domain Consistency and Constraint-Aware Operations}
\label{sec:consistency}
Consistency mandates that a model's outputs adhere to the fundamental physical laws or operational limits inherent to the target domain. Purely data-driven architectures frequently generate unfeasible predictions---such as negative vehicular speeds or resource usage exceeding an operational upper bound. Such violations severely degrade operator trust in critical infrastructure. Inspired by emerging paradigms in safety-critical generation that utilize output-space guidance, a robust forecasting framework must enforce physical feasibility directly within its decoding pathway, moving beyond the limitations of reactive, post-hoc corrections.

\subsubsection{Spatio-Temporal Complexity and Multi-Scale Dynamics}
\label{sec:complexity}
Complexity involves the management of dense spatial couplings across heterogeneous nodes. In physical networks, nodes interact via multiple simultaneous mechanisms: local topological dependencies (e.g., aerodynamic wake effects, adjacent road segments) and global environmental fronts (e.g., synoptic weather patterns). Existing models frequently conflate these multi-scale interactions into a singular spatial operator, thereby failing to distinguish localized perturbations from system-wide drivers.

To systematically resolve this trilemma, we propose PhyST-Net, an ExG-aware Spatio-Temporal Graph Transformer designed as a versatile framework for domain-constrained network forecasting. To accommodate the multi-scale and multi-modal nature of dynamic network signals, we introduce the following key innovations:
\begin{itemize}
    \item \textbf{Continuous and Differentiable Patching}: We present the \textbf{Adaptive Gaussian-soft-windowed Patching (AGSP)} mechanism. By transforming discrete observations into learnable spectral-temporal densities, AGSP recovers continuous event dynamics and functions as an adaptive low-pass filter, effectively rejecting high-frequency noise.
    \item \textbf{Decoupled Spatio-Temporal Topology}: We introduce the \textbf{Relational Spatio-Temporal Manifold Fusion (R-STMF)} framework. R-STMF leverages a dual-branch graph logic to decompose latent node dynamics into multi-scale spatial manifolds, mitigating the spatial conflation typical of standard distance-based graphs.
    \item \textbf{Piecewise Conditioning and Safe Capping}: We implement the \textbf{KAN-enhanced Adaptive Modulation and Conditioning (K-AMC)} module, utilizing Kolmogorov-Arnold spline edge functions for the non-linear modulation of graph-enhanced states via exogenous drivers. This is coupled with a \textbf{Physics-Informed Dynamic Capping Head (PI-DCH)} that embeds strict physical safety constraints and adherence to operational envelopes directly into the output space.
\end{itemize}

By explicitly partitioning the data flow into distinct computational pathways---historical context, relational reasoning, and exogenous conditioning---PhyST-Net establishes a rigorous forecasting contract. This interface-centric design prevents the entanglement of historical data and future drivers, ensuring the architecture can be seamlessly instantiated across diverse critical domains.

\subsection{Contributions}
\label{sec:contributions}
The primary contributions of this work are summarized as follows:
\begin{itemize}
    \item \textbf{A general ExG-aware graph-forecasting formulation.} We define network forecasting as a multi-modal task where historical states, exogenous variables, and domain constraints are processed through distinguishable pathways, thereby supporting broad architectural generalizability.
    \item \textbf{A continuous temporal tokenization module.} We introduce AGSP, which supersedes rigid patch boundaries with learnable Gaussian soft-windows, preserving event continuity and minimizing boundary artifacts in non-stationary sequences.
    \item \textbf{A modular graph adapter for mixed spatial mechanisms.} We develop R-STMF to disentangle multi-scale spatial dependencies prior to relational propagation, providing a modular contract that accommodates alternative graph predictors.
    \item \textbf{A future-ExG-conditioned decoding path.} We fuse K-AMC and PI-DCH to inject forecast-horizon drivers and structurally enforce feasibility constraints within the model, bridging latent representation learning with strict operational requirements.
\end{itemize}

The effectiveness and robustness of PhyST-Net are validated through comprehensive experiments on real-world datasets spanning diverse scales and operational conditions, confirming its viability as a state-of-the-art general-purpose network forecasting framework.
